% LaTeX-ready replacement snippets for manuscript, cover letter, and rebuttal

\section*{Working title}

False-positive-resistant nonlocal diagnosis of field-free topological superconductivity in compensated-magnetic Josephson junctions

\section*{Main-text figure captions}

\paragraph{Figure 1 \textbar{} Diagnostic logic for field-free topology in a compensated-magnetic Josephson junction.}
This figure defines the physical setting and the diagnostic standard used throughout the manuscript. \textbf{a,} Schematic of a phase-biased Josephson junction formed by two s-wave superconducting leads coupled through a compensated-magnetic weak link, with left and right transport probes used to resolve local and nonlocal response. \textbf{b,} Illustration of momentum-dependent spin splitting in the compensated magnetic region: the spin texture changes sign in momentum space while the net magnetization remains zero, distinguishing this mechanism from a conventional Zeeman field. \textbf{c,} Local near-zero-energy features can arise from both topological and trivial states, so a zero-bias or near-zero spectral feature alone is not a sufficient diagnostic. \textbf{d,} Workflow of the combined criterion used here: bulk topological transition, open-boundary spectra, negative-control tests and phase-resolved nonlocal response must be mutually consistent before a near-zero state is identified as topological. The central claim of the study is therefore not the appearance of a near-zero mode by itself, but the closure of this full evidence chain under controlled false-positive tests.

\paragraph{Figure 2 \textbar{} Phase-driven bulk topological transition without net magnetization.}
This figure shows that superconducting phase bias drives a bulk topological transition in the absence of a conventional Zeeman field. \textbf{a,} Normal-state band splitting generated by the compensated magnetic texture, highlighting finite momentum-dependent spin splitting at zero net magnetization. \textbf{b,} Bulk excitation gap as a function of phase bias and chemical potential, revealing a gap closing and reopening that marks the transition between trivial and topological regimes. \textbf{c,} Corresponding $Z_2$ topological invariant over the same parameter space, showing that the gap reopening coincides with a change in bulk topology rather than a crossover between spectrally similar trivial states. \textbf{d,} Representative cut at fixed magnetic texture strength, resolving the gap minimum and reopening across the transition. Together, these panels establish that field-free topology in this junction is controlled by the interplay of compensated magnetic splitting, induced pairing and phase bias, and that the bulk transition can be identified independently of boundary spectroscopy.

\paragraph{Figure 3 \textbar{} Boundary spectra reveal near-zero modes, but zero energy alone does not identify topology.}
This figure separates genuine topological boundary modes from trivial near-zero mimics. \textbf{a,} Open-boundary spectrum versus phase bias, showing the emergence of a near-zero branch only after the bulk transition identified in Fig.~2. \textbf{b,} Spatial profile of the low-energy state in the topological regime, with weight localized at opposite ends of the junction, consistent with a nonlocal boundary mode. \textbf{c,} Counterexample from a trivial parameter regime or local perturbation, where a near-zero state appears but remains spatially asymmetric or insufficiently isolated from the low-energy background. \textbf{d,} Spectral-isolation metric for the lowest mode, demonstrating that the topological branch is distinguished not by zero energy alone, but by simultaneous end localization and clean separation from nearby trivial states. These results motivate the stricter claim used in the main text: local near-zero structure is necessary but not sufficient, and boundary spectroscopy must be interpreted together with bulk topology and nonlocal transport.

\paragraph{Figure 4 \textbar{} Negative controls expose trivial routes to near-zero signatures.}
This figure tests the diagnostic against prespecified trivial null ensembles rather than selected counterexamples. \textbf{a,} Local impurity and weak-disorder perturbations applied in a bulk-trivial regime generate near-zero spectral mimics despite the absence of a topological transition. \textbf{b,} Distribution of the final diagnostic score across topological, impurity-generated trivial and disorder-generated trivial ensembles, showing that local mimics do not reproduce the same ensemble-level behavior as the topological branch. \textbf{c,} Effect-size summary, with bootstrap confidence intervals, quantifying the separation between the topological ensemble and the two null ensembles. \textbf{d,} Empirical true-positive and false-positive rates for the prespecified pass criterion, demonstrating that the trivial null ensembles remain confined to the expected background failure band. \textbf{e,} Stability of the null separation under variation of nuisance parameters such as impurity strength, disorder amplitude or thermal broadening. The figure therefore shows not only that trivial perturbations can mimic local low-energy features, but that they fail the final diagnostic in a statistically calibrated sense.

\paragraph{Figure 5 \textbar{} A single phase-resolved nonlocal criterion identifies the topological window and fails under the same rule in trivial controls.}
This figure converts the transport response into a single prespecified decision system. \textbf{a,} Definition of the raw lead-resolved nonlocal conductance, the antisymmetrized signal $G^{-}(\phi)$, and the fixed threshold-and-window rule used throughout the figure. \textbf{b,} Raw nonlocal response in the clean junction. \textbf{c,} Extraction of the connected phase window selected by the single nonlocal criterion. \textbf{d,} Comparison between that extracted window and the independently determined bulk-topological window. \textbf{e,} Application of the same locked criterion to impurity-driven and disorder-driven trivial controls, which retain local low-energy structure but fail the nonlocal decision rule. \textbf{f,} Final pass-fail summary, including finite-temperature, finite-bias and geometry perturbations under the same criterion. The figure therefore establishes one operational statement only: the topological regime is identified by a reproducible nonlocal decision window that is fixed before control comparisons and is not reproduced by trivial near-zero mimics.

\section*{Fig.~5 reconstruction}

\paragraph{One-line job.}
Fig.~5 must do one thing only: show that the same phase window that carries the bulk topological transition is the only window that produces a robust, lead-resolved nonlocal transport response, while trivial near-zero states fail this transport test.

\paragraph{Recommended six-panel logic.}
\begin{itemize}
\item \textbf{5a | Measurement geometry and observable definition.} Show the probe geometry and define $G_{LR}(V,\phi)$, the antisymmetrized signal $G^{-}(\phi)=G_{LR}(\phi)-G_{LR}(-\phi)$, and the extracted scalar criterion.
\item \textbf{5b | Clean topological map.} Plot $G_{LR}(V,\phi)$ in the clean topological regime with the bulk-transition window marked directly on the panel.
\item \textbf{5c | Aligned phase cut.} Compare, on the same phase axis, the bulk gap, the boundary-state isolation metric and the extracted nonlocal observable.
\item \textbf{5d | Trivial control comparison.} Show that trivial near-zero states do not reproduce the same phase-resolved nonlocal pattern under the same plotting logic.
\item \textbf{5e | Robustness summary.} Test finite temperature, finite bias broadening and representative lead-geometry variation.
\item \textbf{5f | Final operational criterion.} End with a pass-fail summary panel that labels the topological window directly.
\end{itemize}

\paragraph{Expanded caption if Fig.~5 is upgraded to six panels.}
\textbf{Figure 5 \textbar{} A phase-resolved nonlocal transport signal isolates the topological regime from trivial near-zero mimics.} This figure converts the spectral evidence into an operational transport diagnostic. \textbf{a,} Lead geometry and definition of the nonlocal observable used throughout the figure. \textbf{b,} Nonlocal differential conductance $G_{LR}(V,\phi)$ in the clean junction, showing a phase-locked response that emerges in the same phase window as the bulk topological transition. \textbf{c,} Aligned phase cut comparing the bulk gap, boundary-state isolation metric and extracted nonlocal observable, demonstrating that these three quantities select the same topological window. \textbf{d,} Trivial control case, in which local near-zero structure may persist but the phase-resolved nonlocal response is suppressed or loses its topological correlation. \textbf{e,} Robustness of the extracted transport criterion under finite temperature, finite bias averaging and representative lead-geometry variation. \textbf{f,} Final pass-fail summary of the operational criterion, identifying the regime in which nonlocal transport, bulk topology and boundary spectra remain mutually consistent. The figure therefore provides the manuscript's decisive result: topology is not inferred from a local near-zero feature alone, but from a reproducible phase-resolved nonlocal response tied to the same window selected independently by bulk and boundary diagnostics.

\section*{Extended Data and Supplement captions}

\paragraph{Extended Data Figure 1 \textbar{} Robustness of the nonlocal diagnostic window under parameter variation.}
This figure extends the main-text transport analysis beyond the representative parameter set used in Fig.~5. \textbf{a,} Nonlocal response window across a broader chemical-potential range. \textbf{b,} Dependence on compensated magnetic splitting strength. \textbf{c,} Stability against moderate variation of pairing amplitude or phase-bias resolution. \textbf{d,} Summary map of parameter regions in which the combined nonlocal criterion remains valid. The purpose of this figure is not to enlarge the central claim, but to show that the diagnostic window is not an isolated fine-tuned point in parameter space.

\paragraph{Extended Data Figure 2 \textbar{} Finite-size scaling of low-energy boundary modes.}
This figure tests whether the near-zero boundary branch behaves consistently with a nonlocal topological mode under changes in system length. \textbf{a,} Open-boundary low-energy spectrum for increasing junction length. \textbf{b,} Energy splitting of the lowest pair versus system length. \textbf{c,} End-state localization profile at representative lengths. \textbf{d,} Extracted decay trend demonstrating suppression of finite-size hybridization in the topological regime. These data support the interpretation that the topological near-zero branch is controlled by end-to-end hybridization rather than by an accidental local resonance.

\paragraph{Extended Data Figure 3 \textbar{} Lead-geometry dependence of the nonlocal transport signal.}
This figure checks whether the transport criterion survives changes in probe configuration. \textbf{a,} Alternative lead geometries considered in the transport calculations. \textbf{b,} Nonlocal differential conductance for each geometry in the topological regime. \textbf{c,} Comparison of the antisymmetrized nonlocal signal and its phase derivative across geometries. \textbf{d,} Summary of the geometry-dependent diagnostic pass window. The central point is that the precise signal amplitude may vary with contact arrangement, but the topological phase window is still reflected in a reproducible phase-resolved nonlocal response.

\paragraph{Extended Data Figure 4 \textbar{} Thermal and bias broadening of the transport criterion.}
This figure evaluates the experimental visibility of the nonlocal signal beyond the zero-temperature idealization. \textbf{a,} Broadening of the local response with increasing temperature. \textbf{b,} Broadening of the nonlocal response over the same temperature range. \textbf{c,} Effect of finite bias averaging on the extracted criterion. \textbf{d,} Remaining diagnostic window after combined thermal and bias broadening. These results define the regime in which the proposed transport signature should remain experimentally accessible and clarify where the signal becomes too washed out for reliable identification.

\paragraph{Supplementary Figure 1 \textbar{} Lattice model, conventions and numerical workflow.}
This figure documents the computational setup used throughout the manuscript. \textbf{a,} Lattice geometry and boundary conditions. \textbf{b,} Definition of the compensated magnetic texture and superconducting phase convention. \textbf{c,} Schematic of the BdG basis and lead attachment used in bulk, boundary and transport calculations. \textbf{d,} Workflow linking Hamiltonian construction, topological diagnosis, open-boundary diagonalization and transport evaluation. It serves as a reproducibility guide rather than as an additional physics claim.

\paragraph{Supplementary Figure 2 \textbar{} Independent checks of the $Z_2$ topological transition.}
This figure cross-validates the bulk transition reported in Fig.~2 using independent numerical diagnostics. \textbf{a,} Gap closing and reopening along a representative cut. \textbf{b,} Corresponding $Z_2$ invariant. \textbf{c,} Alternative topological indicator or convergence check. \textbf{d,} Agreement of transition boundaries extracted from the different procedures. The figure is included to show that the reported phase boundary is numerically stable and not an artifact of a single diagnostic choice.

\paragraph{Supplementary Figure 3 \textbar{} Open-boundary spectra across representative trivial and topological cuts.}
This figure expands the boundary-state analysis beyond the single main-text example. \textbf{a--c,} Spectra and mode profiles along representative trivial cuts. \textbf{d--f,} Spectra and mode profiles along representative topological cuts. The comparison clarifies which features persist systematically across the phase diagram and which are specific to isolated parameter choices.

\paragraph{Supplementary Figure 4 \textbar{} Ensemble statistics for impurity-induced and disorder-induced false positives.}
This figure provides the full statistical basis for the negative-control argument summarized in Fig.~4. \textbf{a,} Distribution of minimum excitation energies across impurity realizations. \textbf{b,} Distribution of edge-weight asymmetry. \textbf{c,} Disorder-ensemble pass and fail counts for each element of the combined criterion. \textbf{d,} False-positive rate after applying the full diagnostic filter. The figure is essential for demonstrating that the exclusion of trivial mimics is statistical rather than anecdotal.

\paragraph{Supplementary Figure 5 \textbar{} Comparison of candidate transport indicators.}
This figure records the broader indicator search that is intentionally compressed out of the main text. \textbf{a,} Raw nonlocal response. \textbf{b,} Antisymmetrized nonlocal response. \textbf{c,} Phase derivative and related candidate metrics. \textbf{d,} Summary of which indicators remain consistent with the bulk transition and boundary-state window. It justifies the choice of the final criterion while keeping the main manuscript focused on one operational observable.

\paragraph{Supplementary Figure 6 \textbar{} Low-energy effective description near the phase-driven transition.}
This figure connects the numerical results to a compact physical picture. \textbf{a,} Effective low-energy spectrum near the transition point. \textbf{b,} Parameter dependence of the effective mass or gap term. \textbf{c,} Matching between the effective description and the full lattice calculation. \textbf{d,} Limits of validity of the reduced model. The figure is meant to sharpen interpretation, not to enlarge the scope of the claim beyond the numerically established regime.

\paragraph{Supplementary Figure 7 \textbar{} Source-data and code-level reproducibility checks.}
This figure documents the practical reproducibility package accompanying the manuscript. \textbf{a,} Mapping between main-text panels and source-data files. \textbf{b,} Plot-regeneration workflow. \textbf{c,} Numerical convergence and seed bookkeeping. \textbf{d,} File-level checklist for figure reproduction. Its role is to make the numerical claims auditable and easy to reconstruct from the released code and source data.

\section*{Abstracts}

\subsection*{Nature Communications version}
Field-free routes to topological superconductivity are attractive because strong magnetic fields and ferromagnetic stray fields can compromise proximity-induced pairing and complicate spectroscopic interpretation. Yet local near-zero-energy features alone do not distinguish Majorana end modes from trivial Andreev or impurity-induced states. Here we study a phase-biased Josephson junction with a compensated magnetic weak link, in which momentum-dependent spin splitting drives a bulk gap closing and a $Z_2$ topological transition without net magnetization. We show that trivial near-zero states can be generated by local impurities and weak disorder, but they fail a combined diagnostic based on bulk topology, boundary-state isolation and phase-resolved nonlocal transport. The resulting criterion provides a false-positive-resistant route for identifying field-free topological superconductivity in compensated-magnetic superconducting devices.

\subsection*{PRL version}
We show that a compensated-magnetic Josephson junction can host a phase-driven topological superconducting transition without net magnetization and that the transition is tracked by a robust nonlocal transport signature rather than by a local near-zero mode alone. In this system, momentum-dependent spin splitting produces a bulk gap closing and reopening under superconducting phase bias, while open-boundary spectra develop end-localized near-zero states only in the topological regime. By comparing with impurity-induced and disorder-induced trivial near-zero states, we find that local spectroscopy is insufficient but a phase-resolved nonlocal response remains diagnostic. The result identifies a field-free mechanism for topological superconductivity and a transport criterion that sharply distinguishes it from common false positives.

\section*{Introduction}

\paragraph{Paragraph 1.}
Field-free routes to topological superconductivity are becoming increasingly important because the magnetic fields commonly used to drive Majorana platforms can also suppress proximity-induced pairing, introduce orbital complications and blur spectroscopic interpretation. This tension is especially acute in Josephson-based architectures, where phase control offers an attractive tuning knob but where the absence of a clean magnetic-field-free diagnostic has limited how convincingly low-energy boundary features can be interpreted. The field therefore needs not only additional candidate platforms, but also stricter ways to determine when a field-free low-energy state is genuinely topological.

\paragraph{Paragraph 2.}
That need remains unmet. Local zero-bias or near-zero-energy signatures are now understood to be intrinsically ambiguous because trivial Andreev bound states, impurity-induced in-gap states and weak-disorder resonances can all generate similar low-energy features. At the same time, recent planar and compensated-magnetic Josephson-junction proposals have sharpened the platform landscape without resolving the central exclusion problem: how to identify topology without relying on a local near-zero mode alone. As a result, the limiting issue is no longer whether one can engineer another putative field-free junction, but whether one can impose an experimentally relevant false-positive filter on the resulting low-energy states.

\paragraph{Paragraph 3.}
Here we study a phase-biased Josephson junction with a compensated magnetic weak link that produces momentum-dependent spin splitting without net magnetization. In this setting, superconducting phase bias acts as the control parameter that drives a bulk gap closing and reopening together with a $Z_2$ topological transition. This makes the junction a useful arena for separating platform novelty from diagnostic rigor: the compensated magnetic texture supplies a field-free route to the transition, while the phase dependence allows the bulk, boundary and transport responses to be compared on the same footing. The key question is therefore not whether a near-zero mode can appear, but whether all physically relevant observables point to the same topological window.

\paragraph{Paragraph 4.}
We answer that question by constructing a closed evidence chain. We first identify the bulk topological transition from the gap closing and $Z_2$ invariant. We then show that open-boundary near-zero modes become meaningful only when they are spectrally isolated and end-localized in the same phase window. Next, we introduce explicit impurity and disorder controls to test whether trivial low-energy states can mimic the boundary signal. Finally, we show that a phase-resolved nonlocal transport response survives these false-positive tests and provides the operational discriminator missing from local spectroscopy alone. We therefore identify topology not from the appearance of a near-zero mode, but from the consistency of bulk topology, boundary spectra and phase-resolved nonlocal response under controlled false-positive tests.

\paragraph{Aggressive opener option.}
Local near-zero-energy features are no longer the bottleneck in field-free Majorana platforms; credible false-positive exclusion is.

\section*{Results}

\paragraph{Bulk topology is driven by phase bias in the absence of net magnetization.}
We begin by establishing the bulk transition that underlies the rest of the manuscript. In the compensated-magnetic junction, the magnetic texture generates momentum-dependent spin splitting while preserving zero net magnetization, so the relevant control parameter is not a conventional Zeeman field but the superconducting phase bias. Figure~2 shows that varying $\phi$ drives a bulk gap closing and reopening over a finite parameter window, and that this spectral transition is accompanied by a simultaneous change in the $Z_2$ invariant. This alignment is essential: the reopening is not treated as a generic low-energy rearrangement, but as the point at which the bulk diagnosis changes from trivial to topological. The phase-biased compensated-magnetic junction therefore provides a field-free route to a topological transition whose location can be identified directly from bulk quantities before any appeal is made to boundary spectroscopy.

\paragraph{Boundary near-zero modes emerge in the topological window, but zero energy alone remains ambiguous.}
We next turn to the open-boundary spectrum to determine what the bulk transition implies for low-energy boundary structure. As shown in Fig.~3, a near-zero branch appears only after entering the topological phase window identified from the bulk calculation, and representative wavefunction profiles on that side of the transition show weight concentrated at opposite ends of the junction. Taken alone, however, this observation is still insufficient. Trivial parameter choices or local perturbations can also generate low-energy boundary features that approach zero energy without carrying the same topological meaning. For this reason, we do not identify the topological branch by zero energy alone. Instead, the relevant boundary signature is the coexistence of near-zero energy, end localization and spectral isolation within the same phase window selected independently by the bulk transition. Figure~3 therefore plays a narrowing role rather than a concluding one: it shows what a candidate topological boundary mode should look like, while making explicit why local low-energy structure cannot serve as a sufficient criterion.

\paragraph{Negative controls show that trivial low-energy mimics survive local tests but fail the full evidence chain.}
The ambiguity of local spectroscopy becomes most useful when it is converted into a null test rather than a qualitative caution. We therefore evaluate explicit trivial null ensembles, generated in bulk-trivial regimes with either local impurity perturbations or weak disorder, under the same prespecified diagnostic rule used for the topological ensemble. Figure~4 is organized around this comparison. The first point is that the null can indeed generate local near-zero mimics: trivial samples can acquire low-energy structure that would be misleading if judged only by local spectra or selected wavefunction profiles. The second, and more important, point is that these null realizations remain quantitatively separated from the topological ensemble once the full diagnostic score is evaluated at the distribution level.

This separation is not presented as a visual impression alone. Instead, we compare the topological and null ensembles through a fixed scalar diagnostic outcome and report both the associated effect size and the null-calibrated false-positive rate. The relevant question is therefore not whether some trivial samples can look superficially similar, but whether the topological ensemble remains distinguishable from impurity-driven and disorder-driven trivial ensembles after uncertainty is taken into account. Figure~4 shows that it does: the topological ensemble occupies a distinct score distribution, while the null ensembles remain confined to the expected false-positive background band. In this form, the control analysis becomes a quantitative falsification step rather than a gallery of counterexamples.

\paragraph{Phase-resolved nonlocal transport provides the operational discriminator.}
Having established both the bulk transition and the limitations of local boundary spectroscopy, we finally ask which observable can be promoted to a genuine decision rule. Figure~5 is designed to answer that question with a single prespecified nonlocal criterion rather than with a family of related transport indicators. We define the criterion from the antisymmetrized lead-resolved conductance and require that it exceed a fixed threshold over a connected phase interval that overlaps the independently computed bulk-topological window. Under this rule, the role of the transport analysis is no longer to display several suggestive structures, but to decide whether the same nonlocal window appears in the clean topological regime and whether it fails under the same evaluation rule in trivial controls.

This locking step is essential. Without it, the transport section would remain vulnerable to the objection that different observables or thresholds might be chosen retrospectively for different cases. Instead, Fig.~5 applies one criterion to all comparisons. The clean junction produces a reproducible nonlocal decision window that coincides with the bulk and boundary window, whereas impurity-driven and disorder-driven trivial cases retain local low-energy structure but fail the same nonlocal rule. Extended Data Fig.~1 then tests whether this decision window survives finite temperature, finite bias broadening and representative geometry variation without redefining the criterion. The contribution of Fig.~5 is therefore not that many transport quantities change near the transition, but that one locked nonlocal rule yields a topological pass/fail decision that is consistent with the rest of the evidence chain and does not reproduce under trivial near-zero mimics.

\section*{Discussion}

The main implication of these results is not simply that a compensated-magnetic Josephson junction can host field-free topological superconductivity, but that such a platform can be interrogated by a sharper diagnostic standard than local spectroscopy alone. This distinction matters because the bottleneck in the present field is no longer the absence of candidate low-energy platforms. Rather, it is the difficulty of excluding trivial low-energy mimics once local near-zero structure appears. By aligning bulk topology, open-boundary spectra, explicit negative controls and phase-resolved nonlocal transport on the same phase window, the present analysis shifts the focus from platform plausibility to diagnostic closure.

This is also the sense in which the nonlocal transport result should be interpreted. We do not claim that nonlocal response by itself constitutes a universal proof of topology across arbitrary device classes. Instead, within the compensated-magnetic phase-biased junction studied here, the nonlocal signal functions as the operational observable that remains consistent with the independently determined bulk and boundary windows while trivial controls fail the same test. The value of the criterion is therefore comparative and exclusionary: it is stronger than a local near-zero feature because it survives precisely the false-positive scenarios that make local spectroscopy ambiguous.

The results further suggest that field-free Josephson architectures may be especially useful when phase bias can be used as a common axis along which bulk, boundary and transport observables are compared. In that setting, the question is not whether a given low-energy state looks Majorana-like at a single point in parameter space, but whether multiple observables become mutually consistent over a reproducible phase interval. That framing is likely to be more useful for realistic device assessment than a strategy centered on isolated zero-bias features. In this sense, the present work is best viewed as a diagnostic framework for a timely class of junctions, rather than as a claim that the false-positive problem has been solved in full generality.

Several limitations should be stated explicitly. First, the present study is theoretical and numerical, so experimental visibility is inferred through finite-temperature, finite-bias and geometry tests rather than demonstrated in a specific device. Second, the transport criterion is developed for the compensated-magnetic junction class considered here and should not be advertised as model-agnostic without additional cross-platform analysis. Third, the strength of the manuscript depends on the statistical quality of the impurity and disorder controls and on the reproducibility of the source-data and code package. These are not peripheral implementation details; they are part of the evidentiary standard required for the central claim. For that reason, the contribution should be framed as a false-positive-resistant protocol for this junction class, with clearly stated scope and boundary conditions.

Within those limits, the study advances the field in a concrete way. It identifies a field-free route to a bulk topological transition in a compensated-magnetic Josephson junction, shows why local near-zero modes do not close the case, and supplies a lead-resolved nonlocal observable that tracks the same topological window selected independently by bulk and boundary analysis. That combination is the manuscript's real contribution, and it is the level at which the work should be compared with recent Josephson-junction proposals.

\section*{Conclusion}

We have studied a compensated-magnetic Josephson junction in which superconducting phase bias drives a bulk gap closing, reopening and $Z_2$ topological transition without net magnetization. Open-boundary calculations show that near-zero boundary modes emerge in the topological regime, but explicit impurity and disorder controls demonstrate that local low-energy features alone remain insufficient for identification. The key result is that a phase-resolved nonlocal transport response selects the same window identified independently by the bulk transition and isolated boundary spectrum, while trivial near-zero mimics fail the same consistency test. We therefore conclude that, for this class of field-free junctions, topology is most credibly identified not by a local near-zero signature in isolation, but by a false-positive-resistant evidence chain built from bulk topology, boundary-state structure and nonlocal transport.

\section*{Methods / Supplement overview}

The Methods section in the main manuscript should be short, functional and explicitly subordinate to the Results. Its job is not to re-derive the full framework, but to let a specialist understand what was computed, what observables were extracted and where the technical details have been deposited. For this paper, the main-text Methods should do four things only.

First, it should define the model at the level needed to read Figs.~2--5: the compensated-magnetic Josephson-junction geometry, the BdG Hamiltonian, the phase convention, the boundary conditions and the default parameterization used in the main figures. Second, it should state how the three core diagnostics are computed: the bulk gap and $Z_2$ invariant, the open-boundary spectra with end-weight and spectral-isolation metrics, and the lead-resolved nonlocal conductance with the final extracted criterion. Third, it should identify the minimal numerical choices that matter for interpretation, including $k$-grid resolution, energy broadening, disorder sampling scale and finite-temperature convolution. Fourth, it should point the reader explicitly to the Supplement for derivations, convergence checks, alternative indicators, ensemble details and reproducibility assets.

The main-text Methods should remain compact because the manuscript's claim is not methodological novelty in every numerical ingredient. The novelty lies in how these ingredients are combined into a false-positive-resistant evidence chain. Anything that does not directly help the reader understand the origin of the four main Results claims should be moved out of the main text.

\paragraph{What stays in the main-text Methods.}
\begin{itemize}
\item One paragraph defining the lattice/BdG model and compensated magnetic texture.
\item One paragraph explaining how the bulk gap and $Z_2$ invariant are evaluated.
\item One paragraph defining the open-boundary diagnostics: end weight and spectral isolation.
\item One paragraph defining the transport calculation, $G_{LR}$, the antisymmetrized signal and the final operational criterion.
\item A short closing sentence routing detailed derivations, robustness checks and reproducibility information to the Supplement.
\end{itemize}

\paragraph{What must not stay in the main-text Methods.}
\begin{itemize}
\item Full derivations of the topological invariant.
\item Multiple alternative diagnostic indicators that were not kept in the final narrative.
\item Long convergence tables or parameter sweeps.
\item Full impurity/disorder ensemble bookkeeping.
\item Detailed lead self-energy derivations and geometry comparisons.
\item Reproducibility manifests, software inventories or run instructions.
\end{itemize}

\paragraph{Supplement overview.}
The Supplement should not behave like a storage closet. It should be organized as a second evidentiary layer whose sole function is to support, stress-test and delimit the five claims already made in the main manuscript. Every Supplement section should therefore answer one of three questions: how was the main-text quantity defined, how do we know it is numerically stable, or how far does the final diagnostic survive beyond the representative main-text examples?

\paragraph{Supplement S1 | Model definition.}
This section fixes the exact Hamiltonian, basis conventions, geometry, pairing profile, compensated magnetic texture, phase convention and default parameters. Its purpose is to make the main model unambiguous.

\paragraph{Supplement S2 | Topological invariant.}
This section provides the formal definition and numerical implementation of the $Z_2$ invariant, together with convergence checks and at least one independent validation route. Its purpose is to defend Fig.~2 against the attack that the bulk phase boundary is diagnostic-dependent or numerically fragile.

\paragraph{Supplement S3 | Open-boundary calculations.}
This section defines system sizes, diagonalization choices, end-weight measures, spectral-isolation metrics and length scaling. Its purpose is to defend Fig.~3 against the attack that the boundary signal is a finite-size or metric artifact.

\paragraph{Supplement S4 | Impurity and disorder controls.}
This section specifies impurity form, disorder distribution, number of realizations, seed logic, pass/fail criterion and ensemble statistics. Its purpose is to defend Fig.~4 against the attack that the false-positive analysis is anecdotal.

\paragraph{Supplement S5 | Lead-attached transport.}
This section defines lead self-energies, transport observables, EC/CAR-related quantities if shown, energy broadening, finite-temperature convolution and contact-geometry sensitivity. Its purpose is to defend Fig.~5 against the attack that the nonlocal criterion is tuned, fragile or underdefined.

\paragraph{Supplement S6 | Low-energy effective model.}
This section gives the compact interpretive model that explains why the topological branch yields a phase-locked nonlocal response and why impurity-induced trivial states can mimic local peaks without generating the same robust window. Its purpose is interpretive, not claim-expanding.

\paragraph{Supplement S7 | Reproducibility.}
This section enumerates data files, scripts, run commands, software versions, hardware environment, figure-to-file mapping and any source-data packaging notes. Its purpose is to defend the manuscript against reproducibility and transparency objections rather than to add scientific narrative.

\paragraph{Extended Data versus Supplement boundary.}
Extended Data should contain robustness figures that a reviewer may need to inspect visually while judging the main claim: finite-size scaling, parameter-window stability, geometry dependence and thermal/bias broadening. The Supplement should contain definitions, derivations, bookkeeping and expanded statistics. A good rule is this: if a panel helps a reviewer decide whether the main claim survives, it can live in Extended Data; if it explains how the panel was generated or records a fuller parameter archive, it belongs in the Supplement.

\paragraph{Non-negotiable boundary rules.}
\begin{itemize}
\item Main text: one final transport criterion only.
\item Extended Data: robustness of the chosen criterion.
\item Supplement: failed or secondary indicators, not the main narrative.
\item Main text: one representative topological window.
\item Supplement: broader parameter sweeps and convergence evidence.
\item Main text: one clear impurity/disorder message.
\item Supplement: full ensemble details and pass/fail bookkeeping.
\end{itemize}

\paragraph{Short Methods-routing paragraph for the manuscript.}
The system is modeled as a phase-biased compensated-magnetic Josephson junction within a Bogoliubov-de Gennes framework, from which we compute the bulk gap and $Z_2$ invariant, the open-boundary spectra with end-localization metrics, and the lead-resolved nonlocal conductance used in the final diagnostic criterion. The main text reports only the definitions and numerical settings required to interpret Figs.~2--5. Full model specification, invariant implementation, open-boundary diagnostics, disorder and impurity statistics, transport details, effective-model analysis, and reproducibility information are provided in the Supplementary Information.

\section*{Cover letter}

Dear Editor,

We wish to submit the manuscript entitled ``False-positive-resistant nonlocal diagnosis of field-free topological superconductivity in compensated-magnetic Josephson junctions'' for consideration as an Article in \textit{Nature Communications}.

The manuscript addresses a central problem in the present search for Majorana modes: local near-zero-energy features are not, by themselves, reliable evidence for topological superconductivity. At the same time, field-free or low-stray-field routes are increasingly attractive because conventional magnetic-field-based platforms can complicate proximity-induced pairing and blur spectroscopic interpretation. Our work brings these two issues together in a single setting.

We study a phase-biased Josephson junction with a compensated magnetic weak link and show that the system undergoes a bulk gap closing and $Z_2$ topological transition without net magnetization. More importantly, we do not identify topology from a local low-energy feature alone. Instead, we build the evidence chain from four mutually constrained elements: bulk topology, open-boundary spectra, explicit negative controls against impurity- and disorder-induced trivial states, and a phase-resolved nonlocal transport response. This combined logic turns a model study into an operational diagnostic framework for field-free topological superconductivity.

We believe the work fits \textit{Nature Communications} for three reasons. First, it addresses a timely problem of broad interest across topological superconductivity, mesoscopic transport and compensated magnetic systems. Second, it goes beyond proposing another candidate platform by confronting the false-positive problem directly and quantitatively. Third, it emphasizes experimentally relevant observables, especially phase-resolved nonlocal response, rather than relying on local spectroscopy alone.

This manuscript is original, is not under consideration elsewhere, and has been approved by all authors. We hope that you will find it suitable for external review.

Sincerely,\\
[Corresponding author name]\\
On behalf of all authors

\section*{Anticipated editor objections}

\paragraph{Objection 1.}
This looks too close to existing compensated-magnetic or altermagnetic Josephson-junction literature to justify external review.

\paragraph{Recommended framing.}
The manuscript should be pitched not as a new platform announcement, but as a solution to a live diagnostic problem in field-free topological superconductivity. The novelty is the false-positive-resistant nonlocal identification framework, not the bare existence of a compensated-magnetic junction.

\paragraph{Objection 2.}
The central evidence may still rely too heavily on local near-zero structure.

\paragraph{Recommended framing.}
The revised package should state, as early as possible, that local near-zero features are treated as insufficient by design. The decisive evidence is the consistency among bulk topology, boundary-state isolation, negative controls and phase-resolved nonlocal transport.

\paragraph{Objection 3.}
The transport criterion may look tuned to this model rather than broadly useful.

\paragraph{Recommended framing.}
Do not claim universality. Claim an operational criterion for the studied class of compensated-magnetic phase-biased junctions, with experimentally relevant observables and explicit failure tests.

\paragraph{Objection 4.}
The manuscript may not yet be complete enough in robustness analysis for review.

\paragraph{Recommended framing.}
Before submission, the package should make two things unmistakable: the negative controls are statistical rather than illustrative, and the nonlocal signature remains visible after finite-temperature, finite-bias and geometry checks.

\section*{Reviewer response template}

We thank the Editor and the Reviewers for their careful reading of our manuscript and for the constructive comments. The main concern raised across the reports is the strength of the evidence used to distinguish topological boundary modes from trivial near-zero states. In the revised manuscript, we have therefore tightened the central claim, added explicit false-positive controls, clarified the bulk-topology analysis, and strengthened the transport section so that the identification of topology does not rely on local near-zero features alone. Below we respond point by point.

\paragraph{Reviewer \#X, Comment 1.}
[Paste reviewer comment]

\paragraph{Response.}
Thank you for raising this point. We agree that [state the precise issue neutrally]. In the revised manuscript, we addressed this concern in three ways. First, we now [new analysis / new figure / rewritten claim]. Second, we clarified [methodological point, parameter definition, or interpretation boundary] in [section / figure / supplement location]. Third, we have revised the wording throughout the manuscript so that the claim is restricted to [supported claim], rather than the broader formulation used previously. These changes are intended to ensure that [specific criticism] is no longer required as an assumption.

\paragraph{Changes in manuscript.}
\begin{itemize}
\item Main text: [section, paragraph]
\item Figure(s): [figure number / panel]
\item Supplement: [section]
\end{itemize}
